%
% Conclusion.tex
%
% Copyright (C) 2025 by Inspirefly.
%
% 2.305 GHz Circular Patch Antenna Documentation
%

\chapter{Conclusion}

A 2.305~GHz circular microstrip patch antenna for the PocketQube parent board has been sized, simulated, and laid out in a reproducible flow that links analytical cavity theory, MATLAB Antenna Toolbox exploration, ANSYS HFSS full-wave verification, and KiCAD 10.0 fabrication data.

The cavity model (Eqs.~\ref{eq:fres}--\ref{eq:ae}) gives a physical radius of $19.87$~mm on the final FR-408HR stack ($\epsilon_r=3.68$, $h=0.1$~mm prepreg), within $0.5$~mm of the HFSS-tuned radius used in \texttt{CircularPatchV3.aedt}. The early rectangular $\lambda_g/2$ estimate ($L\approx 17.8$~mm on RO4350B) and the hexagonal variant were evaluated and superseded, with provenance retained in git commits \texttt{f49b060}, \texttt{5b98260}, and \texttt{fad1ff9}. The MATLAB particle-swarm search (\texttt{matlab/MPA\_opt\_swrm.m}) demonstrated the limits of a single-layer lossless model, converging to a low-impedance minimum ($Z\approx 0.3$~$\Omega$, RL${}\approx 0$~dB) that confirmed the need for direct HFSS tuning of patch radius and feed offset.

The HFSS model \texttt{CircularPatchV3.aedt} (commit \texttt{235c732}, ``Two coax approach for 1U design'') implements two orthogonal probe feeds for circular polarization. When driven in quadrature through the external APS-02-WBM-LP-DF-CC divider, simulation predicts $S_{11} < -10$~dB per port at 2.305~GHz, inter-port isolation below $-18$~dB, gain near 4--5~dBi, and axial ratio below 3~dB at boresight. Geometry was exported as \texttt{CircularPatchV3.dxf}/\texttt{.step} and converted via \texttt{dxf\_to\_kicad\_mod.py} to the KiCAD footprint \texttt{Inspirefly:CircularPatchAntenna} (commit \texttt{2f72509}). The final board (\texttt{antennaBoard.kicad\_pcb} at \texttt{5554f5c}) is a $51.85 \times 51.89$~mm, 1.3414~mm thick, 4-copper-layer PCB using FR-408HR throughout, with two Molex 732511350 U.FL-compatible coax connectors and through-vias at the tuned feed points. The board fits the 1P face and leaves the PocketQube rail envelope clear.

Future work is verification. A first-article fabrication at the Table~\ref{tab:stackup} stack, followed by VNA measurement of $S_{11}$ and inter-port $S_{21}$, will validate the 50~$\Omega$ match and the $\pm 18$~MHz lot-to-lot frequency sensitivity to $\epsilon_r$. Anechoic or near-field range measurement of realized gain and axial ratio versus angle will confirm the HFSS pattern and the 0.4~dB budget for the external divider and phase-matched cables. If $f_0$ is offset, a single corrective spin of radius $a$ in 0.1~mm steps in HFSS and a corresponding DXF re-export will restore centre frequency without changing the feed network. Long-term, integration with the parent bus and thermal-vacuum testing will verify that the thin prepreg under the patch maintains adhesion and that the U.FL connectors survive vibration.

All sources -- ANSYS AEDT projects in \texttt{ansys/}, MATLAB scripts and logs in \texttt{matlab/}, KiCAD schematics, PCB, footprints, and DXF/STEP exports in \texttt{KiCAD/}, and this document in \texttt{doc/} -- are versioned in \texttt{Pocketqube-parent-antennaBoard} git, with binary KiCAD backups in \texttt{KiCAD/antennaBoard/antennaBoard-backups/} and HFSS-to-KiCAD traceability through commit \texttt{235c732} $\rightarrow$ \texttt{2f72509} $\rightarrow$ \texttt{5554f5c}.

